\chapter*{Jak używać tej książki}\addcontentsline{toc}{chapter}{Jak używać tej książki}\markboth{Jak używać tej książki}{Jak używać tej książki}

Książka ma trzy części i \val{chapters.total} rozdziałów. Część I (rozdziały 1--3) dotyczy czasu,
część II (rozdziały 4--6) wyboru i formatu, a część III (rozdziały 7--10) jakości czytania
i rocznego planu. Rozdziały nie muszą iść w tej kolejności, ale pierwsze trzy są najbardziej
arytmetyczne i warto je przeczytać, zanim ustalisz liczbę książek na rok.

\section*{Tempo: jeden rozdział na tydzień}
Proponuję jeden rozdział tygodniowo. Dziesięć rozdziałów to dziesięć tygodni pracy nad książką,
a pozostałe tygodnie roku to praktyka: czytanie i prowadzenie dziennika. To propozycja, a nie
wymaganie. Jeśli rozdział zajmuje dwa tygodnie, nic się nie stało. Zadanie na tydzień na końcu
każdego rozdziału uznajesz za wykonane wtedy, gdy je zrobiłeś, a nie gdy przeczytałeś rozdział.

\section*{Co jest w każdym rozdziale}
Na początku stoi lista \glqq Po tym rozdziale potrafisz\grqq{}. Sprawdź, czy umiesz to zrobić,
zanim przejdziesz dalej. Tekst korzysta z pięciu rodzajów ramek:
\begin{itemize}
\item \textbf{Przykład} -- rozpisany rachunek z liczbami, który możesz sprawdzić kartką i kalkulatorem;
\item \textbf{Pułapka} -- częsty błąd lub mit wraz z jego korektą;
\item \textbf{Zasada} -- reguła do przyjęcia razem z powodem, dla którego ją podaję;
\item \textbf{Uwaga} -- zastrzeżenie, którego nie chcę pominąć;
\item \textbf{Z półki} -- odniesienie do innej pracy lub praktyki; tam, gdzie znam jej treść, piszę, co w niej jest.
\end{itemize}

\section*{Ćwiczenia i odpowiedzi}
Każdy rozdział kończy się ćwiczeniami. Odpowiedź jest zapisana przy każdym z nich, a wszystkie
odpowiedzi zebrano w dodatku A. Ćwiczenie ma sens tylko wtedy, gdy najpierw spróbujesz je
rozwiązać, dlatego nie zaglądaj do odpowiedzi przed próbą.

\section*{Znaczek [?]}
Pomarańczowy znaczek [?] oznacza twierdzenie o świecie, na przykład o zdrowiu, pamięci, czasie
czytania innych ludzi albo prawie, którego autor nie sprawdził w pierwotnym źródle w tej sesji.
Brak znaczka nie oznacza, że twierdzenie jest prawdziwe. Oznacza tylko, że nie zostało oznaczone
jako niesprawdzone. Liczby policzone w tekście i arytmetyka pokazana krok po kroku znaczka nie
wymagają. Dodatek D opisuje, jak takie twierdzenia sprawdzać.

\section*{Dodatki}
Dodatek A zawiera odpowiedzi do ćwiczeń. Dodatek B podaje tablice liczb oraz wzór, z którego
powstały, wraz z założeniami. Dodatek C zawiera dzienniczki do wydruku: dziennik czytania, siatkę
roku i listę do przeczytania. Dodatek D opisuje źródła i sposób ich weryfikacji. Dodatek Z jest
rejestrem stanu książki: liczby w nim są liczone z danych, a nie wpisywane ręcznie.
